\section{Mapping to GQL and SQL/PGQ}

The path algebra can be used to describe path queries from GQL and SQL/PGQ using a common set of operators \cite{angles2024path}. In particular, the path modes of these languages can be represented using recursive operators together with grouping, ordering, and projection. This makes it possible to describe the evaluation of a path query as a sequence of algebraic operations.

\subsection{Selectors and Restrictors}

GQL uses selectors and restrictors to control the paths produced by a path query \cite{angles2024path}:
\begin{itemize}
    \item \textbf{Restrictors}: determine which paths are considered during recursive evaluation, such as WALK, TRAIL, SIMPLE, and ACYCLIC.
    \item \textbf{Selectors}: determine which of the resulting paths are returned, such as ALL, ANY, and ANY SHORTEST.
\end{itemize}

\subsection{Algebraic Encoding}

Each selector-restrictor combination can be expressed using a combination of:

\begin{itemize}
    \item recursion operator ($\phi$),
    \item grouping operator ($\gamma$),
    \item ordering operator ($\tau$),
    \item projection operator ($\pi$).
\end{itemize}

\subsection{Example Translation}

Consider the following GQL query:

\begin{verbatim}
MATCH ANY SHORTEST TRAIL
p = (x)-[:KNOWS]->+(y)
\end{verbatim}

The query asks for one shortest trail between each pair of connected nodes. The
TRAIL restrictor means that edges cannot be repeated, while ANY SHORTEST
selects one shortest path for each source-target pair \cite{angles2024path}.

The corresponding algebraic evaluation can be understood as a sequence of
operations. First, the edges labeled \texttt{KNOWS} are selected. The recursive
operator $\phi_{\text{Trail}}$ then generates paths that satisfy the TRAIL
semantics. The resulting paths are grouped by their source and target nodes,
ordered by path length, and finally projected so that one path is returned from
each group.

In simplified form, the structure of the translation is:

\[
\pi \Big(
    \tau \big(
        \gamma \big(
            \phi_{\text{Trail}}(
                \sigma_{\text{KNOWS}}(\text{Edges}(G))
            )
        \big)
    \big)
\Big).
\]

The exact parameters of $\gamma$, $\tau$, and $\pi$ determine how paths are
grouped, ordered, and selected. This example shows how a GQL path mode can be
represented as a composition of algebraic operators \cite{angles2024path}.

\subsection{Expressiveness}

The path algebra can represent the selector-restrictor combinations defined by
GQL and SQL/PGQ \cite{angles2024path}. At the same time, the algebra allows
additional combinations because grouping, ordering, projection, and recursion
are represented as separate operators.

For example, the algebra can combine a recursive trail operator with grouping
by path length and then select one path from each group. Such a combination is
not directly provided by the GQL path mode syntax \cite{angles2024path}.

The authors show that the different choices of grouping, ordering, projection,
and recursion result in a larger space of possible expressions than the
selector-restrictor combinations defined by GQL. This means that the algebra
is not limited to reproducing existing language features, but can also describe
path operations that are not directly available in the current standards
\cite{angles2024path}.